% ============================================================================
\clearpage
\section*{v4.9.0-r64 Source-Only: A Frame Audit of the Gravity Statements under the Counting Reading}
\addcontentsline{toc}{section}{v4.9.0-r64 Source-Only: A Frame Audit of the Gravity Statements under the Counting Reading}
\label{sec:v490-r64-gravity-audit}
% ============================================================================

\paragraph{Purpose.}
After v4.9.0-r63 the author asked what gravity is in this context.  Under the counting reading the two values of \(\gamma\) come from the same horizon, with the same area
\(A\) and the same \(\lP\): each counting fixes its own \(\gamma\) by \(S=A/(4\lP^2)\).  This layer audits the main-body statements that use \(\gL\), \(\gBI\) or \(\gstring\) on the
gravity side, as v4.9.0-r61 did for OP1 to OP5.  It makes no prediction and rewrites no earlier layer.

\noindent\textbf{Source-only status.}
This v4.9.0-r64 layer is source-only.  It does not build a PDF, does not assign a new DOI,
and does not authorize a public release.  The v4.9.0 Version DOI
\href{https://doi.org/10.5281/zenodo.23534954}{10.5281/zenodo.23534954} and the Concept DOI
\href{https://doi.org/10.5281/zenodo.19210728}{10.5281/zenodo.19210728} are unchanged by it.

\subsection*{v4.9.0-r64 results}
\label{subsec:v490-r64-results}

\paragraph{Theoremlet v4.9.0-r64.1 (where the frame difference sits).}
\begin{enumerate}[leftmargin=2em]
\item[(a)] Theorem~\ref{thm:k-values} lists \(\gBI\) with \(k=0\), ``frame-invariant by construction''; Theorem~\ref{thm:frame-offset} calls \(\gBI\) the Lorentz-transformed
\(\gstring\).  The two statements contradict each other, independently of r63.
\item[(b)] Under the counting reading the whole frame difference sits in \(\gamma\): \(\hbar\), \(G_N\) and \(\lP\) enter both countings identically.  The assignments
\(k(\hbar)=1\), \(k(G_N)=2\) converted with \(\gL\) use the counting ratio a second time.  This concerns Theorem~\ref{thm:planck-photon} (\(\gL^{-3/2}\)), the corollary
``Cosmological constant revisited'' (\(\gL^{5}\)) and item~(i) of Proposition~\ref{prop:hawking-gBI-tachyon} (\(\gL^{-1}\) for \(T_H\)).  They are not refuted, but they need a
source for their \(k\) values other than \(\gL\).
\item[(c)] Theorem~\ref{thm:frame-offset} assigns \(\gstring\) to the photon frame and reads \(\gL\) as an M--R Lorentz factor at \(\beta=0.844\).  The number is \(\tanh2\eta\), the
M--T velocity; under r63 \(\gstring\) belongs to the tachyon floor.  The value is consistent, the frame label is R for T.
\item[(d)] Theorem~\ref{thm:symmetric-frame} is frame-consistent and gains a kinematic reading from r63.4: the tachyon at \(-\bsym\) is the floor of the mirror frame.
\item[(e)] Theorem~\ref{thm:hawking-triolic} gives the escaping photon in R the value \(\gstring\).  Under r63 that is the counting of the infalling tachyon; the statement
reads T as R.
\item[(f)] Theorem~\ref{thm:dual-origin} (\(\gL\approx1/\bsym\), gap \(2.4\%\)) compares \(\cosh2\eta\) with \(\coth\eta\), the class mixing of r61.1(a).  Under r63 \(\gL\) is a
counting ratio; the gap is not geometric.
\item[(g)] Theorem~\ref{thm:dark-energy} uses \(\bsym^2=\tanh^2\eta\), the M--R relation, which is consistent within the symmetric placement; its value \(\bsym=1/\sqrt3\) is the
Hopf value of r61.1(b).  The drift problem of v4.9.0-r17 is untouched.
\item[(h)] Theorem~\ref{thm:einstein} uses none of the three quantities.
\end{enumerate}

\paragraph{Proof sketch.}
(a) by reading the two statements; (b) from the counting equations of r63.1, in which \(\lP\) and \(A\) are common to both countings; (c)--(g) by comparison with r63.4,
v4.9.0-r18/r19 and r61; (h) by inspection.  \(\square\)

\paragraph{Reading.}
\begin{enumerate}[leftmargin=2em]
\item The gravity statements split into two groups.  Those that put the frame difference into \(\gamma\) (Theorems~\ref{thm:frame-offset} after relabelling,
\ref{thm:symmetric-frame}) are compatible with the counting reading.  Those that put it into \(\hbar\) and \(G_N\) (Theorems~\ref{thm:k-values}, \ref{thm:planck-photon}, the
cosmological corollary, Proposition~\ref{prop:hawking-gBI-tachyon}(i)) are not supported by it.
\item In the counting reading gravity itself is frame-independent: \(G_N\), \(\lP\) and the area are the same for all frames.  What changes between the frames is how many
states a frame counts per area, and \(\gamma\) is the exchange rate.
\item This is a reading, not a theorem: the \(k\)-value picture would survive with an independent derivation of its exponents.  Candidates for status notes are
Theorems~\ref{thm:k-values} (r64.1(a)), \ref{thm:planck-photon} and \ref{thm:hawking-triolic}; the decision is left to the author.
\item H15 (\(\cot\theta_W=\gL\)) is not affected.
\end{enumerate}

\subsection*{v4.9.0-r64 status in the three frames}
\label{subsec:v490-r64-frames}

\begin{center}
\small
\begin{tabular}{@{}>{\raggedright\arraybackslash}p{0.30\textwidth}>{\raggedright\arraybackslash}p{0.62\textwidth}@{}}
\toprule
result & M / T / P\\
\midrule
r64.1(a), (b) & frame difference in the counting (M, T), not in \(\hbar\), \(G_N\)\\
r64.1(c), (e) & T read as R\\
r64.1(d) & M and T; frame-consistent\\
r64.1(f), (g) & as r61.1(a), (b)\\
\bottomrule
\end{tabular}
\end{center}

\subsection*{v4.9.0-r64 register entries}
\label{subsec:v490-r64-register}

\begin{center}
\small
\begin{tabular}{@{}>{\raggedright\arraybackslash}p{0.34\textwidth}>{\raggedright\arraybackslash}p{0.58\textwidth}@{}}
\toprule
statement & status after v4.9.0-r64\\
\midrule
\(k(\gBI)=0\) (Theorem~\ref{thm:k-values}) & contradicts Theorem~\ref{thm:frame-offset}; candidate for a status note\\
photon-frame Planck length, \(\gL^5\) vacuum ratio, Hawking (i) & not supported by the counting reading; need independent \(k\) values\\
Hawking temperature \(\gstring T_H\) & reads the tachyon counting as the photon value\\
symmetric matter--tachyon frame & frame-consistent; T is the floor of the mirror frame\\
\bottomrule
\end{tabular}
\end{center}

\subsection*{v4.9.0-r64 ledger and guard update}
\label{subsec:v490-r64-ledger-guard}

\[
R_{\rm grav}^{v490\text{-}r64}=(1,1,0)
\]
for (contradiction in the \(k\) table found; frame difference located in the counting under r63; no gravity statement derived).

The layer spends no guard:
\[
R_{\rm glob}^{(v490\text{-}r64,+)}=(0,1,1,1,1),
\qquad
\Gamma_{\rm guard}^{v490\text{-}r64}=(1,1,1,1,1),
\qquad
\beta_4=1/256\ \text{unspent}.
\]
It does not imply exact finite \(C5\) zero closure, exact all-readout terminality,
\(C6/T2\), a global photon-clock/RH theorem, or global all-readout closure.
